% ID: SEL-D0030-OA-E04
% ESTADO: validada. S031, cotejada con examen y criterios.
\textbf{a)} Tomamos dos tornillos distintos, sin reemplazo,
y describimos cada muestra sin orden. Las $\binom42=6$
muestras equiprobables y sus medias en gramos son
\[
\begin{array}{c|rrrrrr}
\text{Muestra}&(1,2)&(1,3)&(1,4)&(2,3)&(2,4)&(3,4)\\\hline
\bar x&1.5&2&2.5&2.5&3&3.5
\end{array}
\]

\textbf{b)} La esperanza de la media muestral es
$(1.5+2+2.5+2.5+3+3.5)/6=\boxed{2.5\text{ g}}$.
Promediamos las desviaciones cuadráticas con peso $1/6$,
pues calculamos la varianza de una distribución completa:
\[
\operatorname{Var}(\bar X)=\frac{1+0.25+0+0+0.25+1}{6}
=\boxed{5/12\text{ g}^2}.
\]

% AUD001: alternativa para hipótesis no explícita.

\textit{Hipótesis de muestreo.} El enunciado no explicita la reposición.
Si se distingue el orden sin reemplazo, las doce parejas equiprobables
producen la misma distribución de medias que los seis subconjuntos.

Con reemplazo, en cambio, hay dieciséis pares ordenados equiprobables.
En la tabla, la fila representa la primera extracción y la columna
la segunda. Cada celda da la media del par:
\[
\begin{array}{c|cccc}
i\backslash j&1&2&3&4\\\hline
1&1&1{,}5&2&2{,}5\\
2&1{,}5&2&2{,}5&3\\
3&2&2{,}5&3&3{,}5\\
4&2{,}5&3&3{,}5&4
\end{array}
\]
La esperanza sigue siendo $2{,}5$. La varianza poblacional es
$[(1-2{,}5)^2+(2-2{,}5)^2+(3-2{,}5)^2+(4-2{,}5)^2]/4=5/4$.
Por independencia de las dos extracciones,
$\operatorname{Var}(\bar X)=(5/4)/2=\boxed{5/8}$.
La elección entre $5/12$ y $5/8$ depende de la hipótesis de reposición.
